
\documentclass[a4paper,14pt]{extarticle}
\usepackage{fontspec}
\usepackage{polyglossia}
\setmainlanguage[babelshorthands=true]{russian}
\setotherlanguage{english}
%\setmainfont{Tinos}
%\setsansfont{DejaVu Sans}
%\setmonofont{DejaVu Sans Mono}

\setmainfont{Times New Roman}
\setsansfont{Arial}
\setmonofont{Courier New}

%\newfontfamily\cyrillicfont{Tinos}
%\newfontfamily\cyrillicfonttt{DejaVu Sans Mono}

\newfontfamily\cyrillicfont{Times New Roman}
\newfontfamily\cyrillicfonttt{Courier New}

\usepackage{amsmath,amssymb}
\usepackage{geometry}
\geometry{left=30mm,right=15mm,top=20mm,bottom=20mm}
\usepackage{setspace}
\onehalfspacing
\usepackage{indentfirst}
\setlength{\parindent}{1.25cm}
\usepackage{array,booktabs,longtable,tabularx}
\usepackage{listings}
\usepackage{xcolor}
\usepackage{caption}
\usepackage{hyperref}
\hypersetup{hidelinks}
\renewcommand{\lstlistingname}{Листинг}
\renewcommand{\lstlistlistingname}{Список листингов}
\lstset{
  basicstyle=\ttfamily\footnotesize,
  breaklines=true,
  columns=fullflexible,
  keepspaces=true,
  frame=single,
  numbers=left,
  numberstyle=\tiny,
  tabsize=2,
  captionpos=b
}
\captionsetup[table]{justification=centering,labelsep=colon}
\captionsetup[lstlisting]{justification=centering,labelsep=colon}
\begin{document}

\begin{titlepage}
\fontsize{12}{14}\selectfont
\thispagestyle{empty}
\begin{center}
МИНИСТЕРСТВО НАУКИ И ВЫСШЕГО ОБРАЗОВАНИЯ РОССИЙСКОЙ ФЕДЕРАЦИИ

\vspace{4pt}
{\normalsize федеральное государственное автономное образовательное учреждение высшего образования “Санкт-Петербургский политехнический университет Петра Великого”

\vspace{4pt}
Институт компьютерных наук и кибербезопасности

\vspace{4pt}
Высшая школа технологий искусственного интеллекта

\vspace{4pt}
02.03.01 Математика и компьютерные науки}

\vspace{1.9cm}
{\large Отчет по курсовой работе

\vspace{4pt}
{\normalsize Разработка игрового бота для расширенной версии игры “Крестики-нолики”}

\vspace{4pt}
{\normalsize по дисциплине “Структуры данных”}}
\end{center}

\vspace{2.2cm}
\begin{flushleft}
Выполнил:

Студент группы 5130201/50302\hspace{3.0cm}\rule{3.5cm}{0.5pt}\hspace{3pt}Кузнецов М.С.

\vspace{10pt}
Проверил:

Преподаватель\hspace{6.0cm}\rule{3.5cm}{0.5pt}\hspace{8pt}Сеннов Н.В.
\end{flushleft}

\vspace{0.7cm}
\begin{flushright}
“\rule{1cm}{0.5pt}”\rule{2.1cm}{0.5pt}\hspace{3pt}2026 г.
\end{flushright}

\vfill
\begin{center}
Санкт-Петербург, 2026
\end{center}
\end{titlepage}

\setcounter{page}{2}
\tableofcontents

\clearpage
\section*{Введение}
\addcontentsline{toc}{section}{Введение}

В данной работе рассматривается разработка игрового бота для расширенной версии игры “Крестики-нолики”. В отличие от классической игры на поле $3 \times 3$, рассматриваемая версия проводится на поле $20 \times 20$, может содержать препятствия, а победной считается непрерывная линия из пяти одинаковых символов по горизонтали, вертикали или диагонали. Из-за большого размера поля простой полный перебор продолжений быстро становится непрактичным, поэтому при проектировании игрока использовались эвристические методы, основанные на оценке локальных и позиционных признаков. Такой подход соответствует общим принципам построения алгоритмов для дискретных задач и методам принятия решений в игровых задачах искусственного интеллекта [1, 2].

Предметная область работы относится к классу комбинаторных игр на клеточном поле. В таких играх качество решения определяется не только формальной проверкой победной линии, но и способностью алгоритма заранее оценивать угрозы, перспективные направления развития атаки и влияние занятых клеток на дальнейшие ходы [3]. Для рассматриваемой игры дополнительно важно асимметричное правило завершения партии: игрок $O$ выигрывает сразу после построения линии, а игрок $X$ после своей линии должен учитывать последний ответный ход игрока $O$. Эта особенность меняет оценку позиции и делает стратегию для двух знаков различной.

Цель работы состоит в разработке бота MyPlayer, который по текущему состоянию поля выбирает допустимую клетку для хода за ограниченное время. Для достижения этой цели в работе формализована модель поля, описано множество рассматриваемых ходов, введены оценки атаки и защиты, реализована обработка немедленных побед и блокировок, а также проведены серии автоматических игр против базовых игроков. Исходный код был написан на C++, а при описании реализации использовались справочные материалы по языку [4]. При редактировании пояснительного текста применялся инструмент ChatGPT\footnote{OpenAI. ChatGPT: программный инструмент на основе большой языковой модели [Электронный ресурс]. URL: \url{https://openai.com/chatgpt/} (дата обращения: 10.05.2026).}.

\clearpage
\section{Постановка задачи}

Необходимо реализовать игрока для уже заданной игровой среды расширенных “Крестиков-ноликов”. Логика игрового движка, смена ходов, хранение истории партии и определение общего результата игры не разрабатываются заново; бот получает состояние партии через предоставленный интерфейс и должен вернуть только один следующий ход. Таким образом, задача состоит не в реализации всей игры, а в разработке алгоритма выбора клетки для текущего игрока.

На вход алгоритму подаются размеры поля $N$ и $M$, матрица состояния клеток $A$, знак текущего игрока $s \in \{X,O\}$, а также информация о клетках, недоступных для хода. В данной работе эксперименты проводятся для поля $20 \times 20$. Каждая клетка может быть свободной, занятой символом $X$, занятой символом $O$ или заблокированной препятствием. Ход допустим только в свободную клетку внутри границ поля. На выходе алгоритм должен выдать координаты одной допустимой клетки $(i^*,j^*)$, в которую игровой движок поставит символ текущего игрока.

Алгоритм должен учитывать пять групп условий. Во-первых, он обязан возвращать только допустимый ход. Во-вторых, при наличии немедленной победы текущего игрока ход должен выбираться с учётом правила последнего ответа для $X$. В-третьих, при отсутствии собственной немедленной победы необходимо блокировать ход соперника, который приводит к его победной линии. В-четвёртых, обычные ходы должны сравниваться по численной оценке, включающей атакующие возможности, защиту, близость к активной области и влияние препятствий. В-пятых, алгоритм должен работать достаточно быстро для серийного запуска партий, поэтому перебор всех игровых продолжений не используется.

К сдаче предъявляются количественные требования по скорости и силе игры. Алгоритм должен быть эффективным: время вычисления одного хода должно составлять величину порядка $50$ "--- $100$ мс на рабочей станции или ноутбуке. Алгоритм также должен быть достаточно сильным: процент побед над игроком базового уровня сложности без учёта ничьих должен быть больше $50\%$, а над игроком повышенного уровня сложности "--- больше $25\%$. Оба базовых игрока предоставляются как подключаемая библиотека с закрытым исходным кодом, которая подключается к проекту; соответствующий релиз библиотеки доступен по адресу: \url{https://github.com/vsennov/tictactoe-course/releases/tag/v1.1.0}.

Результатом работы является программная реализация бота MyPlayer и экспериментальная проверка его поведения. Корректность решения оценивается по допустимости возвращаемых ходов, выполнению указанных требований к скорости и силе игры, а также по среднему времени выбора хода.

\clearpage
\section{Математическое описание алгоритма}

\subsection{Модель поля и допустимого хода}

Игровое поле рассматривается как матрица
\begin{equation}
A = (a_{ij}), \qquad i=1,2,\ldots,N, \quad j=1,2,\ldots,M,
\label{eq:board}
\end{equation}
где $a_{ij}$ обозначает состояние клетки в строке $i$ и столбце $j$. Множество возможных состояний клетки задаётся как
\begin{equation}
S = \{\varnothing, X, O, B\},
\label{eq:states}
\end{equation}
где $\varnothing$ означает свободную клетку, $X$ и $O$ "--- символы игроков, а $B$ "--- препятствие. Ход $(i,j)$ допустим тогда и только тогда, когда $1 \leq i \leq N$, $1 \leq j \leq M$ и $a_{ij}=\varnothing$.

Для проверки линий используются четыре направления
\begin{equation}
D = \{(1,0),(0,1),(1,1),(1,-1)\}.
\label{eq:directions}
\end{equation}
Если для некоторого направления $(d_i,d_j) \in D$ пять последовательных клеток
\[
(i,j),\ (i+d_i,j+d_j),\ (i+2d_i,j+2d_j),\ (i+3d_i,j+3d_j),\ (i+4d_i,j+4d_j)
\]
находятся внутри поля и содержат один и тот же знак игрока, то этот игрок построил победную линию длины $5$.

\subsection{Правило победы для игроков X и O}

Пусть $L_s(A)$ означает наличие на поле $A$ линии из пяти символов игрока $s$. Для игрока $O$ победа определяется прямым условием
\begin{equation}
\mathrm{Win}_O(A) \Longleftrightarrow L_O(A).
\label{eq:win_o}
\end{equation}
Для игрока $X$ учитывается последний ответ игрока $O$. Обозначим через $R_O(A)$ количество свободных клеток, ход в которые позволяет игроку $O$ немедленно построить линию из пяти символов. Тогда гарантированная победа игрока $X$ задаётся условием
\begin{equation}
\mathrm{Win}_X(A) \Longleftrightarrow L_X(A) \wedge R_O(A)=0.
\label{eq:win_x}
\end{equation}
Если после хода $X$ линия $X$ уже построена, но $R_O(A)>0$, такой ход рассматривается не как гарантированная победа, а как позиция, в которой $O$ может последним ответом добиться ничьей. Поэтому в оценочной функции для $X$ защитная составляющая имеет больший вес, чем для $O$.

\subsection{Свободные области и центр кластера}

Из-за препятствий геометрический центр поля не всегда является хорошей областью для первого хода и для позиционного выбора. Поэтому вводится понятие свободного кластера. Свободный кластер "--- это максимальное множество свободных клеток, в котором любые две клетки связаны цепочкой соседних свободных клеток. Соседними считаются клетки, имеющие общую сторону или вершину, то есть используется связность по восьми направлениям.

Пусть $K=\{(i_1,j_1),\ldots,(i_k,j_k)\}$ "--- крупнейший свободный кластер. Его средняя точка определяется формулами
\begin{equation}
\bar{i}=\frac{1}{k}\sum_{t=1}^{k} i_t, \qquad
\bar{j}=\frac{1}{k}\sum_{t=1}^{k} j_t.
\label{eq:cluster_mean}
\end{equation}
Центром кластера считается такая свободная клетка $(i_c,j_c)\in K$, для которой евклидово расстояние до точки $(\bar{i},\bar{j})$ минимально. Если несколько клеток имеют одинаковое расстояние, выбирается клетка с большим количеством потенциальных отрезков длины $5$ и меньшим штрафом за близость к препятствиям.

\subsection{Множество ходов-кандидатов}

Пусть $F$ "--- множество всех свободных клеток поля:
\[
F = \{(i,j) \mid a_{ij}=\varnothing\}.
\]
После начала партии алгоритм не оценивает все клетки из $F$, потому что удалённые от уже поставленных символов клетки редко влияют на ближайшую тактику. Множество кандидатов $C$ определяется как
\begin{equation}
C = \{(i,j)\in F \mid \exists (p,q): a_{pq}\in\{X,O\},\ |i-p|\leq 2,\ |j-q|\leq 2\}.
\label{eq:candidates}
\end{equation}
Если на поле ещё нет символов игроков, то используется специальное правило первого хода. Если множество $C$ оказалось пустым из-за изолированного расположения свободных клеток, алгоритм возвращается к рассмотрению всех клеток из $F$.

\subsection{Первый ход}

Для первого хода нельзя использовать близость к уже поставленным символам, поэтому каждая свободная клетка оценивается по трём параметрам. Пусть $L(i,j)$ "--- количество отрезков длины $5$, которые проходят через клетку $(i,j)$ и полностью состоят из свободных клеток; $P(i,j)$ "--- штраф за препятствия в окрестности радиуса $2$; $Q(i,j)$ "--- манхэттенское расстояние до центра крупнейшего свободного кластера. Тогда оценка первого хода равна
\begin{equation}
F_0(i,j)=1000L(i,j)-20P(i,j)-Q(i,j).
\label{eq:first_score}
\end{equation}
Первым ходом выбирается клетка
\begin{equation}
(i^*,j^*)=\arg\max_{(i,j)\in F} F_0(i,j).
\label{eq:first_argmax}
\end{equation}
Такая формула задаёт все используемые веса явно: потенциальная линия получает вес $1000$, близость к препятствию штрафуется с коэффициентом $20$, а расстояние до центра кластера используется как слабый позиционный критерий.

\subsection{Оценка отрезков длины 5}

Для оценки атаки рассматриваются все отрезки длины $5$ в направлениях из формулы \eqref{eq:directions}. Для выбранного игрока $s$ в каждом отрезке подсчитываются три величины: $k$ "--- количество символов игрока $s$, $e$ "--- количество свободных клеток, $b$ "--- количество заблокированных клеток. К заблокированным относятся символы соперника и препятствия. Оценка одного отрезка равна
\[
\mathrm{Score}(k,e,b)=
\begin{cases}
0, & b>0,\\
100000000, & k\geq 5,\\
3000000, & k=4,\ e=1,\\
80000, & k=3,\ e=2,\\
3000, & k=2,\ e=3,\\
30, & k=1,\ e=4,\\
0, & \text{иначе.}
\end{cases}
\]
Атакующая ценность позиции для игрока $s$ определяется как сумма оценок по всем подходящим отрезкам:
\[
A(A,s)=\sum \mathrm{Score}(k,e,b).
\]

\subsection{Тактические угрозы и итоговая оценка хода}

Пусть $A'$ "--- состояние поля после предполагаемого хода текущего игрока $s$ в клетку $(i,j)$, а $\bar{s}$ "--- знак соперника. Через $R_s(A')$ обозначается количество ходов, после которых игрок $s$ может немедленно построить линию из пяти символов. Тактическая часть оценки задаётся формулами
\begin{equation}
T_{\mathrm{danger}}(i,j)=-50000000\cdot R_{\bar{s}}(A'),
\label{eq:danger}
\end{equation}
\[
T_{\mathrm{attack}}(i,j)=
\begin{cases}
30000000, & R_s(A')\geq 2,\\
3000000, & R_s(A')=1,\\
0, & R_s(A')=0,
\end{cases}
\]
\[
T(i,j)=T_{\mathrm{attack}}(i,j)+T_{\mathrm{danger}}(i,j).
\]

Позиционная близость к активной зоне учитывается функцией $N(i,j)$. В окрестности радиуса $2$ собственный символ на манхэттенском расстоянии $1$ даёт $30$ баллов, собственный символ на расстоянии $2$ даёт $8$ баллов, символ соперника на расстоянии $1$ даёт $15$ баллов, а символ соперника на расстоянии $2$ даёт $4$ балла. Близость к центру свободного кластера задаётся функцией
\[
Q_c(i,j)=-(|i-i_c|+|j-j_c|).
\]

Итоговая оценка различается для игроков $X$ и $O$:
\begin{equation}
F_O(i,j)=T(i,j)+22A(A',O)-24A(A',X)+N(i,j)+Q_c(i,j),
\label{eq:o_score}
\end{equation}
\begin{equation}
F_X(i,j)=T(i,j)+12A(A',X)-24A(A',O)+N(i,j)+Q_c(i,j).
\label{eq:x_score}
\end{equation}
Игрок $O$ получает более высокий атакующий коэффициент, поскольку его линия немедленно завершает партию. Игрок $X$ сильнее ориентирован на защиту, так как после его линии возможен последний ответ $O$.

\subsection{Итоговый алгоритм выбора хода}

На вход алгоритму подаются матрица $A$, размеры поля $N$ и $M$, знак текущего игрока $s$ и множество состояний клеток из формулы \eqref{eq:states}. На выходе алгоритм выдаёт координаты допустимого хода $(i^*,j^*)$. Порядок выбора хода следующий.

\begin{enumerate}
\item Найти крупнейший свободный кластер и его центр по формуле \eqref{eq:cluster_mean}.
\item Если на поле нет символов $X$ и $O$, выбрать первый ход по формулам \eqref{eq:first_score} и \eqref{eq:first_argmax} и завершить алгоритм.
\item Сформировать множество ходов-кандидатов $C$ по формуле \eqref{eq:candidates}.
\item Для каждого хода из $C$ проверить, создаёт ли он немедленную победу текущего игрока. Для $O$ используется условие \eqref{eq:win_o}; для $X$ используется условие \eqref{eq:win_x}. Если такой ход существует, выбрать его и завершить алгоритм.
\item Если собственной немедленной победы нет, проверить, существует ли ход соперника, который на следующем ходе создаёт его победную линию. Если такая клетка найдена, занять её текущим ходом.
\item Если срочного хода нет, для каждой клетки из $C$ построить временное состояние $A'$ и вычислить оценку по формуле \eqref{eq:o_score} для $O$ или по формуле \eqref{eq:x_score} для $X$.
\item Вернуть клетку с максимальной итоговой оценкой:
\[
(i^*,j^*)=\arg\max_{(i,j)\in C} F_s(i,j).
\]
\end{enumerate}

\subsection{Обоснование выбранного подхода}

Полный перебор дерева игры на поле $20 \times 20$ приводит к слишком большому числу продолжений, поэтому в работе используется эвристический алгоритм. Он оценивает не всю будущую партию, а признаки, которые непосредственно влияют на ближайшие ходы: готовые линии, открытые отрезки, немедленные угрозы, двойные угрозы, положение относительно активной зоны и влияние препятствий. Такой подход не гарантирует оптимальную игру во всех позициях, но позволяет получать ход за малое время и сохранять разумное качество решений в длинных сериях партий.

\clearpage
\section{Особенности реализации}

Реализация выполнена на языке C++ в классе MyPlayer. Класс содержит два поля: $m\_sign$, в котором хранится знак текущего игрока, и $m\_name$, в котором хранится имя экземпляра игрока. Метод set\_sign получает знак, назначенный игровым движком, метод get\_name возвращает имя игрока, а метод make\_move принимает состояние партии и возвращает координаты следующего хода. Интерфейс класса приведён в приложении~\ref{app:hpp}, а полный код реализации приведён в приложении~\ref{app:cpp}.

Внутри алгоритма состояние поля преобразуется во вспомогательную структуру Board. Она хранит число столбцов, число строк и последовательность значений клеток. Математически эта структура соответствует матрице $A$ из формулы \eqref{eq:board}; одномерное хранение используется только как способ размещения данных в памяти. Доступ к клетке осуществляется по её координатам, а проверки допустимости хода сводятся к проверке принадлежности координат границам поля и проверке равенства клетки состоянию $\varnothing$.

Вспомогательные функции разделены по смыслу. Функции inside, is\_free и is\_player\_sign проверяют границы поля, свободность клетки и принадлежность значения одному из игроков. Функция get\_candidate\_cells строит множество $C$ из формулы \eqref{eq:candidates}. Функции has\_line\_from\_cell и has\_line\_after\_move проверяют наличие победной линии в четырёх направлениях из формулы \eqref{eq:directions}. Функции count\_immediate\_wins и count\_immediate\_wins\_limited вычисляют количество немедленных угроз $R_s(A)$, используемое в формулах \eqref{eq:win_x} и \eqref{eq:danger}.

Поиск крупнейшего свободного кластера реализован обходом в ширину. Для каждой найденной компоненты свободных клеток вычисляется размер и средняя точка, после чего выбирается свободная клетка, ближайшая к этой точке. При равенстве расстояний учитывается число возможных отрезков длины $5$ и штраф за препятствия. Этот порядок соответствует математическому описанию из подраздела 2.3 и не зависит от особенностей C++.

Основной метод make\_move реализует порядок, описанный в подразделе 2.8. Сначала он обрабатывает первый ход, затем ищет собственную немедленную победу, затем блокирует немедленную победу соперника, после чего оценивает оставшиеся кандидаты. Для игрока $X$ отдельно сохраняется ход, который строит линию $X$, но позволяет $O$ последним ходом добиться ничьей; такой ход используется как запасной вариант, если обычная оценка позиции оказывается неблагоприятной.

Реализация является детерминированной: случайный выбор не используется, поэтому при одинаковом состоянии поля бот возвращает один и тот же ход. Это упрощает отладку и делает результаты серийных экспериментов воспроизводимыми.

\clearpage
\section{Результаты работы программы}

Для проверки бота были проведены автоматические серии игр в стандартном игровом движке проекта. Во всех сериях использовалось поле $20 \times 20$ с теми же правилами победы и препятствий, для которых разрабатывался алгоритм. Для каждого внешнего соперника запускались две серии по $1000$ партий: в первой серии MyPlayer играл за $X$, во второй "--- за $O$. Дополнительно была проведена серия из $1000$ партий MyPlayer против MyPlayer. Такая постановка эксперимента позволяет отдельно оценить влияние знака игрока, поскольку правила для $X$ и $O$ несимметричны.

В таблице~\ref{tab:results} приведены абсолютные и процентные результаты игр. Проценты рассчитаны относительно $1000$ партий в каждой серии.


\begin{table}[htbp]
\centering
\caption{Результаты серийного тестирования бота MyPlayer}
\label{tab:results}
\scriptsize
\setlength{\tabcolsep}{2.2pt}
\renewcommand{\arraystretch}{1.18}
\begin{tabularx}{\textwidth}{>{\raggedright\arraybackslash}X >{\centering\arraybackslash}p{1.25cm} >{\centering\arraybackslash}p{1.7cm} >{\centering\arraybackslash}p{1.75cm} >{\centering\arraybackslash}p{1.45cm} >{\centering\arraybackslash}p{1.75cm}}
\toprule
Серия & Знак & Победы MyPlayer & Победы соперника & Ничьи & Победы без ничьих \\
\midrule
MyPlayer "--- BaselineEasy & $X$ & 819 / 81,9\% & 92 / 9,2\% & 89 / 8,9\% & 89,9\% \\
BaselineEasy "--- MyPlayer & $O$ & 556 / 55,6\% & 334 / 33,4\% & 110 / 11,0\% & 62,5\% \\
MyPlayer "--- BaselineHard & $X$ & 185 / 18,5\% & 232 / 23,2\% & 583 / 58,3\% & 44,4\% \\
BaselineHard "--- MyPlayer & $O$ & 156 / 15,6\% & 502 / 50,2\% & 342 / 34,2\% & 23,7\% \\
MyPlayer "--- MyPlayer & $X/O$ & 877 / 87,7\% & 119 / 11,9\% & 4 / 0,4\% & 88,1\% \\
\bottomrule
\end{tabularx}
\end{table}

Результаты из таблицы~\ref{tab:results} показывают, что против BaselineEasy разработанный бот имеет устойчивое преимущество. В сумме за две серии против этого соперника MyPlayer выиграл $1375$ партий из $2000$, то есть $68{,}75\%$ всех партий. Если исключить ничьи, то на $1801$ результативную партию приходится $1375$ побед MyPlayer, что составляет $76{,}3\%$ и превышает требуемые $50\%$. Против BaselineHard результат заметно слабее: всего было получено $341$ победа MyPlayer, $734$ победы соперника и $925$ ничьих. Без учёта ничьих доля побед MyPlayer равна $31{,}7\%$, поэтому формальное требование $25\%$ также выполнено, хотя при игре за $O$ результат хуже, чем при игре за $X$. Это говорит о том, что локальные эвристики хорошо справляются с простым соперником, но против более сильной стратегии им не всегда хватает поиска на несколько ходов вперёд.

Необычный результат наблюдается в серии MyPlayer против MyPlayer: первый экземпляр бота выиграл $87{,}7\%$ партий. Вероятное объяснение состоит в преимуществе первого хода и в агрессивном развитии инициативы после выбора активной области. Поскольку оба экземпляра используют одинаковую детерминированную стратегию, первый игрок раньше занимает центральные перспективные клетки и чаще первым создаёт угрозы, которые второму игроку приходится блокировать.

Время выбора хода измерялось средствами игрового движка. Работа была выполнена на ноутбуке Asus с процессором 13th Gen Intel(R) Core(TM) i7-1355U 1.70 GHz, оперативной памятью 32,00 ГБ и 64-разрядной операционной системой. Средние значения приведены в таблице~\ref{tab:time}.

\begin{table}[htbp]
\centering
\caption{Среднее время расчёта хода и обработки партии}
\label{tab:time}
\footnotesize
\setlength{\tabcolsep}{3pt}
\renewcommand{\arraystretch}{1.2}
\begin{tabularx}{\textwidth}{>{\raggedright\arraybackslash}X >{\centering\arraybackslash}p{2.25cm} >{\centering\arraybackslash}p{2.25cm} >{\centering\arraybackslash}p{2.35cm}}
\toprule
Серия & Ход MyPlayer, мс & Ход соперника, мс & Игровой процесс, мс \\
\midrule
MyPlayer "--- BaselineEasy & 3,4036 & 0,0025 & 1,7132 \\
BaselineEasy "--- MyPlayer & 3,3977 & 0,0025 & 1,7097 \\
MyPlayer "--- BaselineHard & 4,5028 & 0,0028 & 2,2596 \\
BaselineHard "--- MyPlayer & 4,5816 & 0,0029 & 2,3002 \\
MyPlayer "--- MyPlayer & 1,0336 / 1,1970 & "--- & 1,1185 \\
\bottomrule
\end{tabularx}
\end{table}

Данные из таблицы~\ref{tab:time} показывают, что даже при проверке немедленных угроз, обходе кандидатов и вычислении оценок отрезков длины $5$ среднее время выбора хода остаётся в пределах нескольких миллисекунд. Это значительно меньше требуемого диапазона $50$ "--- $100$ мс на рабочей станции или ноутбуке. Время против BaselineEasy и BaselineHard оказалось больше, чем в режиме MyPlayer против MyPlayer, потому что партии против базовых игроков чаще уходили в глубокую стадию, когда почти всё поле уже занято и приходится анализировать больше сложных локальных конфигураций. В самоигре наблюдалась полная доминация первого игрока $X$, из-за чего партии завершались быстрее, а среднее число оцениваемых поздних позиций было меньше.

\clearpage
\section*{Заключение}
\addcontentsline{toc}{section}{Заключение}

В ходе работы был разработан бот MyPlayer для расширенной версии игры “Крестики-нолики” на поле $20 \times 20$. Основным результатом стала детерминированная эвристическая стратегия, которая принимает на вход состояние поля и знак текущего игрока, строит множество ходов-кандидатов, проверяет немедленную победу и необходимость блокировки, а затем выбирает ход по оценочной функции. В алгоритме явно учтены препятствия, центр крупнейшего свободного кластера, потенциальные отрезки длины $5$, двойные угрозы и асимметричное правило последнего хода для игрока $X$.

Практическая проверка показала, что бот уверенно превосходит простого соперника BaselineEasy: суммарно он выиграл $1375$ партий из $2000$, что составляет $68{,}75\%$ всех партий и $76{,}3\%$ результативных партий без учёта ничьих. Против BaselineHard суммарная доля побед без учёта ничьих составила $31{,}7\%$, поэтому формальное требование к игроку повышенного уровня сложности также выполнено. Среднее время выбора хода составило примерно $3{,}4$ мс против BaselineEasy и $4{,}5$"---$4{,}6$ мс против BaselineHard, что значительно меньше требуемого времени порядка $50$ "--- $100$ мс на ноутбуке.

При выполнении работы были закреплены навыки формализации игровой задачи, проектирования эвристической функции, работы с матричным представлением поля, реализации обхода в ширину для поиска связных областей, обработки специальных правил игры и анализа результатов статистического тестирования. Главным ограничением решения остаётся отсутствие полноценного поиска в глубину: бот хорошо оценивает ближайшие угрозы, но может не видеть сложные многоходовые комбинации. Дальнейшее развитие алгоритма целесообразно связать с minimax-поиском по ограниченному множеству лучших кандидатов, автоматическим подбором весов и более точным различением открытых и закрытых линий.

\clearpage
\section*{Список литературы}
\addcontentsline{toc}{section}{Список литературы}
\begin{enumerate}
\item Кормен Т.~Х., Лейзерсон Ч.~Е., Ривест Р.~Л., Штайн К. Алгоритмы: построение и анализ. 3-е изд. М.: Вильямс, 2013. 1328 с.
\item Рассел С., Норвиг П. Искусственный интеллект: современный подход. 4-е изд. М.: Диалектика, 2021. 1152 с.
\item Beck J. Combinatorial Games: Tic-Tac-Toe Theory. Cambridge: Cambridge University Press, 2008. 732 p.
\item C++ reference: справочник по языку C++ [Электронный ресурс]. URL: \url{https://en.cppreference.com/w/} (дата обращения: 10.05.2026).
\end{enumerate}

\clearpage
\appendix
\section{Исходный код заголовочного файла my\_player.hpp}
\label{app:hpp}

В приложении приведён полный заголовочный файл класса MyPlayer. Этот файл содержит объявление полей класса, конструктора и методов, которые используются игровым движком. Основная роль интерфейса описана в разделе 3, а сам интерфейс приведён в листинге~\ref{lst:hpp}.

\begin{lstlisting}[language=C++,caption={Файл my\_player.hpp},label={lst:hpp}]
#pragma once
#include "core/game.hpp"
namespace ttt::my_player {
using game::Event;
using game::IPlayer;
using game::Point;
using game::Sign;
using game::State;
class MyPlayer : public IPlayer {
Sign m_sign = Sign::NONE;
const char *m_name;
public:
MyPlayer(const char *name) : m_sign(Sign::NONE), m_name(name)
{}
void set_sign(Sign sign) override;
Point make_move(const State &game) override;
const char *get_name() const override;
};
}; // namespace ttt::my_player

\end{lstlisting}

\clearpage
\section{Исходный код файла реализации my\_player.cpp}
\label{app:cpp}

В приложении приведён полный файл реализации разработанного бота. В нём находятся функции подготовки поля, поиска кандидатов, проверки линий, оценки позиции и выбора итогового хода. Подробное описание основных групп функций дано в разделе 3, а полный код реализации помещён в приложении~\ref{app:cpp}.

\begin{lstlisting}[language=C++,caption={Файл my\_player.cpp},label={lst:cpp}]
#include "my_player.hpp"
#include <cmath>
#include <limits>
#include <queue>
#include <vector>
namespace ttt::my_player {
static const int WIN_LENGTH = 5;
void MyPlayer::set_sign(Sign sign) { m_sign = sign; }
const char *MyPlayer::get_name() const { return m_name; }
struct Board {
int cols;
int rows;
std::vector<Sign> cells;
Sign get(int x, int y) const {
return cells[y * cols + x];
}
void set(int x, int y, Sign sign) {
cells[y * cols + x] = sign;
}
};
struct ClusterInfo {
bool valid;
int center_x;
int center_y;
int size;
};
static Point make_point(int x, int y);
static bool inside(const Board &board, int x, int y);
static bool is_free(const Board &board, int x, int y);
static bool is_player_sign(Sign sign);
static int count_possible_lines_through_cell(
const Board &board,
int x,
int y
);
static int obstacle_penalty_near_cell(
const Board &board,
int x,
int y
);
static ClusterInfo find_largest_free_cluster_center(const Board
&board);
static Sign get_enemy(Sign sign) {
if (sign == Sign::X)
return Sign::O;
if (sign == Sign::O)
return Sign::X;
return Sign::NONE;
}
static bool is_player_sign(Sign sign) {
return sign == Sign::X || sign == Sign::O;
}
static bool inside(const Board &board, int x, int y) {
return x >= 0 && x < board.cols && y >= 0 && y < board.rows;
}
static bool is_free(const Board &board, int x, int y) {
return inside(board, x, y) && board.get(x, y) == Sign::NONE;
}
static Point make_point(int x, int y) {
Point p;
p.x = x;
p.y = y;
return p;
}
static Board make_board(const State &state) {
Board board;
board.cols = state.get_opts().cols;
board.rows = state.get_opts().rows;
board.cells.resize(board.cols * board.rows);
for (int y = 0; y < board.rows; ++y) {
for (int x = 0; x < board.cols; ++x) {
board.set(x, y, state.get_value(x, y));
}
}
return board;
}
static bool board_empty(const Board &board) {
for (int y = 0; y < board.rows; ++y) {
for (int x = 0; x < board.cols; ++x) {
if (is_player_sign(board.get(x, y))) {
return false;
}
}
}
return true;
}
static int count_free_cells(const Board &board) {
int result = 0;
for (int y = 0; y < board.rows; ++y) {
for (int x = 0; x < board.cols; ++x) {
if (board.get(x, y) == Sign::NONE) {
++result;
}
}
}
return result;
}
static bool has_neighbor(const Board &board, int x, int y, int
radius) {
for (int dy = -radius; dy <= radius; ++dy) {
for (int dx = -radius; dx <= radius; ++dx) {
if (dx == 0 && dy == 0)
continue;
int nx = x + dx;
int ny = y + dy;
if (!inside(board, nx, ny))
continue;
if (is_player_sign(board.get(nx, ny))) {
return true;
}
}
}
return false;
}
static std::vector<Point> get_candidate_cells(const Board &board
) {
std::vector<Point> result;
bool empty = board_empty(board);
for (int y = 0; y < board.rows; ++y) {
for (int x = 0; x < board.cols; ++x) {
if (!is_free(board, x, y))
continue;
if (empty || has_neighbor(board, x, y, 2)) {
result.push_back(make_point(x, y));
}
}
}
if (!result.empty()) {
return result;
}
for (int y = 0; y < board.rows; ++y) {
for (int x = 0; x < board.cols; ++x) {
if (is_free(board, x, y)) {
result.push_back(make_point(x, y));
}
}
}
return result;
}
static bool has_line_from_cell(
const Board &board,
int x,
int y,
Sign sign
) {
const int directions[4][2] = {
{1, 0},
{0, 1},
{1, 1},
{1, -1}
};
for (int d = 0; d < 4; ++d) {
int dx = directions[d][0];
int dy = directions[d][1];
int count = 1;
int nx = x + dx;
int ny = y + dy;
while (inside(board, nx, ny) && board.get(nx, ny) == sign) {
++count;
nx += dx;
ny += dy;
}
nx = x - dx;
ny = y - dy;
while (inside(board, nx, ny) && board.get(nx, ny) == sign) {
++count;
nx -= dx;
ny -= dy;
}
if (count >= WIN_LENGTH) {
return true;
}
}
return false;
}
static bool has_line_after_move(
const Board &board,
Point move,
Sign sign
) {
Board copy = board;
copy.set(move.x, move.y, sign);
return has_line_from_cell(copy, move.x, move.y, sign);
}
static int count_immediate_wins(
const Board &board,
Sign sign
) {
std::vector<Point> cells = get_candidate_cells(board);
int result = 0;
for (Point move : cells) {
if (has_line_after_move(board, move, sign)) {
++result;
}
}
return result;
}
static int count_immediate_wins_limited(
const Board &board,
Sign sign,
int limit
) {
std::vector<Point> cells = get_candidate_cells(board);
int result = 0;
for (Point move : cells) {
if (has_line_after_move(board, move, sign)) {
++result;
if (result >= limit) {
return result;
}
}
}
return result;
}
// X поставил линию.
// Это реальная победа X только если O не может последним ходом
собрать линию.
static bool is_real_x_win_after_move(
const Board &board,
Point move
) {
Board after_x = board;
after_x.set(move.x, move.y, Sign::X);
if (!has_line_from_cell(after_x, move.x, move.y, Sign::X)) {
return false;
}
// Если поле заполнено, O не может сделать последний ход.
if (count_free_cells(after_x) == 0) {
return true;
}
int o_answers = count_immediate_wins_limited(after_x, Sign::O,
1);
return o_answers == 0;
}
// X поставил линию, но O может ответить своей линией.
// Это не победа X, а гарантированная ничья.
static bool is_x_draw_after_move(
const Board &board,
Point move
) {
Board after_x = board;
after_x.set(move.x, move.y, Sign::X);
if (!has_line_from_cell(after_x, move.x, move.y, Sign::X)) {
return false;
}
if (count_free_cells(after_x) == 0) {
return false;
}
int o_answers = count_immediate_wins_limited(after_x, Sign::O,
1);
return o_answers > 0;
}
static long long pattern_score(
int own_count,
int empty_count,
int blocked_count
) {
if (blocked_count > 0) {
return 0;
}
if (own_count >= 5) {
return 100000000LL;
}
if (own_count == 4 && empty_count == 1) {
return 3000000LL;
}
if (own_count == 3 && empty_count == 2) {
return 80000LL;
}
if (own_count == 2 && empty_count == 3) {
return 3000LL;
}
if (own_count == 1 && empty_count == 4) {
return 30LL;
}
return 0;
}
static long long evaluate_board_for_sign(
const Board &board,
Sign sign
) {
const int directions[4][2] = {
{1, 0},
{0, 1},
{1, 1},
{1, -1}
};
Sign enemy = get_enemy(sign);
long long score = 0;
for (int y = 0; y < board.rows; ++y) {
for (int x = 0; x < board.cols; ++x) {
for (int d = 0; d < 4; ++d) {
int dx = directions[d][0];
int dy = directions[d][1];
int end_x = x + (WIN_LENGTH - 1) * dx;
int end_y = y + (WIN_LENGTH - 1) * dy;
if (!inside(board, end_x, end_y))
continue;
int own_count = 0;
int empty_count = 0;
int blocked_count = 0;
for (int i = 0; i < WIN_LENGTH; ++i) {
int nx = x + i * dx;
int ny = y + i * dy;
Sign cell = board.get(nx, ny);
if (cell == sign) {
++own_count;
} else if (cell == Sign::NONE) {
++empty_count;
} else if (cell == enemy) {
++blocked_count;
} else {
++blocked_count;
}
}
score += pattern_score(own_count, empty_count,
blocked_count);
}
}
}
return score;
}
static int neighbor_bonus(
const Board &board,
Point move,
Sign my_sign
) {
int bonus = 0;
Sign enemy = get_enemy(my_sign);
for (int dy = -2; dy <= 2; ++dy) {
for (int dx = -2; dx <= 2; ++dx) {
if (dx == 0 && dy == 0)
continue;
int x = move.x + dx;
int y = move.y + dy;
if (!inside(board, x, y))
continue;
Sign cell = board.get(x, y);
int distance = std::abs(dx) + std::abs(dy);
int value = distance <= 1 ? 30 : 8;
if (cell == my_sign) {
bonus += value;
} else if (cell == enemy) {
bonus += value / 2;
}
}
}
return bonus;
}
static int strategic_center_bonus(
Point move,
const ClusterInfo &cluster
) {
if (!cluster.valid) {
return 0;
}
int distance =
std::abs(move.x - cluster.center_x) +
std::abs(move.y - cluster.center_y);
return -distance;
}
// Специальный выбор атакующего хода за O.
// Ищем ход, после которого у O появляется одна или несколько
// немедленных победных угроз на следующий ход.
// При этом не выбираем ход, после которого X сразу получает
линию.
static bool find_o_attack_move(
const Board &board,
const std::vector<Point> &cells,
const ClusterInfo &cluster,
Point &best_move
) {
bool found = false;
long long best_score = std::numeric_limits<long long>::min();
for (Point move : cells) {
Board after_o = board;
after_o.set(move.x, move.y, Sign::O);
// Если после нашего хода X сразу может собрать линию,
// такой ход слишком опасен.
int x_wins = count_immediate_wins_limited(after_o, Sign::X,
1);
if (x_wins > 0) {
continue;
}
int o_wins = count_immediate_wins_limited(after_o, Sign::O,
2);
if (o_wins == 0) {
continue;
}
long long score = 0;
// Две и более угрозы почти всегда сильнее одной.
if (o_wins >= 2) {
score += 150000000LL;
} else {
score += 30000000LL;
}
score += evaluate_board_for_sign(after_o, Sign::O) * 12;
score -= evaluate_board_for_sign(after_o, Sign::X) * 8;
score += neighbor_bonus(board, move, Sign::O);
score += strategic_center_bonus(move, cluster);
if (!found || score > best_score) {
found = true;
best_score = score;
best_move = move;
}
}
return found;
}
static long long evaluate_move(
const Board &board,
Point move,
Sign my_sign,
const ClusterInfo &cluster
) {
Sign enemy = get_enemy(my_sign);
Board after_my_move = board;
after_my_move.set(move.x, move.y, my_sign);
// Особое правило:
// O выигрывает сразу, X выигрывает только если O не может
ответить линией.
if (my_sign == Sign::O) {
if (has_line_from_cell(after_my_move, move.x, move.y, Sign::
O)) {
return 1000000000LL;
}
}
if (my_sign == Sign::X) {
if (has_line_from_cell(after_my_move, move.x, move.y, Sign::
X)) {
if (count_free_cells(after_my_move) == 0) {
return 1000000000LL;
}
int o_answers = count_immediate_wins_limited(after_my_move
, Sign::O, 1);
if (o_answers == 0) {
return 1000000000LL;
}
// Это не победа, а ничья. Хорошо как fallback,
// но не считаем это абсолютной победой.
return 10000000LL;
}
}
long long score = 0;
int enemy_immediate_wins = count_immediate_wins(after_my_move,
enemy);
int my_immediate_wins_next = count_immediate_wins(
after_my_move, my_sign);
if (enemy_immediate_wins > 0) {
if (my_sign == Sign::O) {
score -= 120000000LL * enemy_immediate_wins;
} else {
score -= 50000000LL * enemy_immediate_wins;
}
}
if (my_sign == Sign::O) {
if (my_immediate_wins_next >= 2) {
score += 150000000LL;
} else if (my_immediate_wins_next == 1) {
score += 30000000LL;
}
} else {
if (my_immediate_wins_next >= 2) {
score += 30000000LL;
} else if (my_immediate_wins_next == 1) {
score += 3000000LL;
}
}
long long attack = evaluate_board_for_sign(after_my_move,
my_sign);
long long defense = evaluate_board_for_sign(after_my_move,
enemy);
if (my_sign == Sign::O) {
// O должен атаковать, но не может игнорировать угрозы X.
// При слабой защите O слишком часто проигрывает.
score += attack * 22;
score -= defense * 24;
} else {
// X должен защищаться сильнее, потому что после линии X у O
есть последний ход.
score += attack * 12;
score -= defense * 24;
}
score += neighbor_bonus(board, move, my_sign);
score += strategic_center_bonus(move, cluster);
return score;
}
static int count_possible_lines_through_cell(
const Board &board,
int x,
int y
) {
const int directions[4][2] = {
{1, 0},
{0, 1},
{1, 1},
{1, -1}
};
int result = 0;
for (int d = 0; d < 4; ++d) {
int dx = directions[d][0];
int dy = directions[d][1];
for (int shift = -(WIN_LENGTH - 1); shift <= 0; ++shift) {
int start_x = x + shift * dx;
int start_y = y + shift * dy;
bool valid = true;
for (int i = 0; i < WIN_LENGTH; ++i) {
int nx = start_x + i * dx;
int ny = start_y + i * dy;
if (!inside(board, nx, ny) || board.get(nx, ny) != Sign
::NONE) {
valid = false;
break;
}
}
if (valid) {
++result;
}
}
}
return result;
}
static int obstacle_penalty_near_cell(
const Board &board,
int x,
int y
) {
int penalty = 0;
for (int dy = -2; dy <= 2; ++dy) {
for (int dx = -2; dx <= 2; ++dx) {
if (dx == 0 && dy == 0)
continue;
int nx = x + dx;
int ny = y + dy;
if (!inside(board, nx, ny))
continue;
Sign cell = board.get(nx, ny);
if (cell != Sign::NONE && cell != Sign::X && cell != Sign
::O) {
int distance = std::abs(dx) + std::abs(dy);
if (distance <= 1) {
penalty += 20;
} else {
penalty += 7;
}
}
}
}
return penalty;
}
static ClusterInfo find_largest_free_cluster_center(const Board
&board) {
ClusterInfo best_cluster;
best_cluster.valid = false;
best_cluster.center_x = 0;
best_cluster.center_y = 0;
best_cluster.size = 0;
std::vector<int> visited(board.cols * board.rows, 0);
const int directions[8][2] = {
{1, 0},
{-1, 0},
{0, 1},
{0, -1},
{1, 1},
{1, -1},
{-1, 1},
{-1, -1}
};
for (int y = 0; y < board.rows; ++y) {
for (int x = 0; x < board.cols; ++x) {
if (!is_free(board, x, y)) {
continue;
}
int start_index = y * board.cols + x;
if (visited[start_index]) {
continue;
}
std::queue<Point> queue;
std::vector<Point> component;
visited[start_index] = 1;
queue.push(make_point(x, y));
long long sum_x = 0;
long long sum_y = 0;
while (!queue.empty()) {
Point current = queue.front();
queue.pop();
component.push_back(current);
sum_x += current.x;
sum_y += current.y;
for (int d = 0; d < 8; ++d) {
int nx = current.x + directions[d][0];
int ny = current.y + directions[d][1];
if (!inside(board, nx, ny)) {
continue;
}
int index = ny * board.cols + nx;
if (visited[index]) {
continue;
}
if (!is_free(board, nx, ny)) {
continue;
}
visited[index] = 1;
queue.push(make_point(nx, ny));
}
}
int component_size = static_cast<int>(component.size());
if (component_size == 0) {
continue;
}
Point center_cell = component.front();
long long best_distance = std::numeric_limits<long long>::
max();
long long best_cell_score = std::numeric_limits<long long
>::min();
for (Point cell : component) {
long long dx = static_cast<long long>(cell.x) *
component_size - sum_x;
long long dy = static_cast<long long>(cell.y) *
component_size - sum_y;
long long distance = dx * dx + dy * dy;
int possible_lines = count_possible_lines_through_cell(
board,
cell.x,
cell.y
);
int obstacle_penalty = obstacle_penalty_near_cell(
board,
cell.x,
cell.y
);
long long cell_score = 0;
cell_score += possible_lines * 1000LL;
cell_score -= obstacle_penalty * 20LL;
if (distance < best_distance ||
(distance == best_distance && cell_score >
best_cell_score)) {
best_distance = distance;
best_cell_score = cell_score;
center_cell = cell;
}
}
if (!best_cluster.valid || component_size > best_cluster.
size) {
best_cluster.valid = true;
best_cluster.center_x = center_cell.x;
best_cluster.center_y = center_cell.y;
best_cluster.size = component_size;
}
}
}
return best_cluster;
}
static Point choose_first_move(
const Board &board,
const ClusterInfo &cluster
) {
Point best = make_point(0, 0);
long long best_score = std::numeric_limits<long long>::min();
for (int y = 0; y < board.rows; ++y) {
for (int x = 0; x < board.cols; ++x) {
if (!is_free(board, x, y))
continue;
int possible_lines = count_possible_lines_through_cell(
board, x, y);
if (possible_lines == 0)
continue;
int obstacle_penalty = obstacle_penalty_near_cell(board, x
, y);
int distance_to_cluster_center = 0;
if (cluster.valid) {
distance_to_cluster_center =
std::abs(x - cluster.center_x) +
std::abs(y - cluster.center_y);
}
long long score = 0;
score += possible_lines * 1000LL;
score -= obstacle_penalty * 20LL;
score -= distance_to_cluster_center;
if (score > best_score) {
best_score = score;
best = make_point(x, y);
}
}
}
return best;
}
Point MyPlayer::make_move(const State &state) {
Board board = make_board(state);
ClusterInfo free_cluster = find_largest_free_cluster_center(
board);
std::vector<Point> cells = get_candidate_cells(board);
if (cells.empty()) {
if (free_cluster.valid) {
return make_point(free_cluster.center_x, free_cluster.
center_y);
}
return make_point(0, 0);
}
if (board_empty(board)) {
return choose_first_move(board, free_cluster);
}
Point draw_by_x = cells.front();
bool has_draw_by_x = false;
// 1. Немедленная победа с учетом особого правила X/O.
if (m_sign == Sign::O) {
// O выигрывает сразу.
for (Point move : cells) {
if (has_line_after_move(board, move, Sign::O)) {
return move;
}
}
} else if (m_sign == Sign::X) {
// X выигрывает только если O не сможет ответить своей
линией.
for (Point move : cells) {
if (is_real_x_win_after_move(board, move)) {
return move;
}
}
// Если есть ход X в линию, но O может ответить линией,
// это гарантированная ничья. Запоминаем как fallback.
for (Point move : cells) {
if (is_x_draw_after_move(board, move)) {
draw_by_x = move;
has_draw_by_x = true;
break;
}
}
}
// 2. Блок немедленной победы соперника.
// Для X это критично: O выигрывает сразу.
// Для O тоже блокируем линию X, потому что иначе X получает
сильнейший темп,
// а O часто не успевает реализовать ответ.
if (m_sign == Sign::X) {
for (Point move : cells) {
if (has_line_after_move(board, move, Sign::O)) {
return move;
}
}
} else if (m_sign == Sign::O) {
// Сначала обязательно блокируем немедленную линию X.
for (Point move : cells) {
if (has_line_after_move(board, move, Sign::X)) {
return move;
}
}
// Если срочного блока нет, ищем активный атакующий ход за O
.
Point o_attack_move = cells.front();
if (find_o_attack_move(board, cells, free_cluster,
o_attack_move)) {
return o_attack_move;
}
}
Point best_move = cells.front();
long long best_score = std::numeric_limits<long long>::min();
for (Point move : cells) {
long long score = evaluate_move(board, move, m_sign,
free_cluster);
if (score > best_score) {
best_score = score;
best_move = move;
}
}
// Если мы играем за X и обычная эвристика видит плохую
позицию,
// но есть ход, который гарантирует ничью, выбираем ничью.
if (m_sign == Sign::X && has_draw_by_x && best_score < 0) {
return draw_by_x;
}
return best_move;
}
} // namespace ttt::my_player
46__

\end{lstlisting}

\end{document}
